\documentclass[12pt]{article}
\usepackage[spanish]{babel}
\usepackage{amsmath, amssymb}
\usepackage{geometry}
\usepackage{listings}
\usepackage{xcolor}
\geometry{margin=1in}

\title{Aplicaciones y Ejercicios de Probabilidad\(Con validación en Python)}
\author{}
\date{}

\lstset{
language=Python,
basicstyle=\ttfamily\small,
keywordstyle=\color{blue},
commentstyle=\color{gray},
stringstyle=\color{red},
breaklines=true
}

\begin{document}

\maketitle

\section{T3. Aplicaciones de Probabilidad Conjunta}

\subsection{Aplicación 1: Caso Discreto}

Se define la función de probabilidad conjunta:

[
P(X=x, Y=y) =
\begin{cases}
0.1 & (x,y)=(0,0) \
0.2 & (x,y)=(0,1) \
0.3 & (x,y)=(1,0) \
0.4 & (x,y)=(1,1)
\end{cases}
]

Validación: la suma total debe ser 1.

\begin{lstlisting}
import numpy as np

prob = {
(0,0):0.1,
(0,1):0.2,
(1,0):0.3,
(1,1):0.4
}

print(sum(prob.values()))
\end{lstlisting}

\subsection{Aplicación 2: Caso Continuo}

Se define la función de densidad:

[
f(x,y) = 2 \quad \text{para } 0 < x < y < 1
]

Validación:

[
\int_0^1 \int_0^y 2 , dx , dy = 1
]

\begin{lstlisting}
import sympy as sp

x, y = sp.symbols('x y')
f = 2

resultado = sp.integrate(sp.integrate(f, (x, 0, y)), (y, 0, 1))
print(resultado)
\end{lstlisting}

\section{T4. Distribuciones de Probabilidad}

\subsection{Distribución Binomial}

\textbf{Ejercicio 1:}

[
P(X=3)=\binom{5}{3}(0.5)^3(0.5)^2
]

\begin{lstlisting}
from math import comb

n=5
p=0.5
k=3

prob = comb(n,k)*(p**k)*((1-p)**(n-k))
print(prob)
\end{lstlisting}

\textbf{Ejercicio 2:}

[
\mu = np = 5(0.5) = 2.5
]

\subsection{Distribución Poisson}

\textbf{Ejercicio 1:}

[
P(X=2)= e^{-3} \frac{3^2}{2!}
]

\begin{lstlisting}
import math

lam=3
k=2

prob = math.exp(-lam)*(lam**k)/math.factorial(k)
print(prob)
\end{lstlisting}

\textbf{Ejercicio 2:}

[
\mu = \lambda = 3
]

\subsection{Distribución Normal}

\textbf{Ejercicio 1:}

[
P(X < 1)
]

\begin{lstlisting}
from scipy.stats import norm

prob = norm.cdf(1, loc=0, scale=1)
print(prob)
\end{lstlisting}

\textbf{Ejercicio 2:}

[
\mu, \quad \sigma^2
]

\subsection{Distribución Exponencial}

\textbf{Ejercicio 1:}

[
P(X < 2)
]

\begin{lstlisting}
from scipy.stats import expon

prob = expon.cdf(2, scale=1)
print(prob)
\end{lstlisting}

\textbf{Ejercicio 2:}

[
\mu = \frac{1}{\lambda}
]

\section{Conclusión}

Se desarrollaron aplicaciones de probabilidad conjunta para casos discreto y continuo, así como ejercicios teóricos de distribuciones importantes, validados mediante Python.

\end{document}
